% =====================================================================
%  starter-book.tex — copy this to start a new BOOK.
%
%  A book has a cover, a transliteration chart, a table of contents and
%  numbered chapters.
%
%  BUILD — from the kit's main folder:   bash tools/build.sh starter-book
%  (LuaLaTeX, three passes: the contents, the cover's page anchors, and
%  one to settle.)
%
%  THE ONE RULE THAT IS NOT TYPOGRAPHIC
%    Nothing here inserts content you did not type — no honorific, no
%    grading, no translation, no ḥadīth number. If you want it on the
%    page, type it.
% =====================================================================
\documentclass[12pt]{article}
\newcommand{\ibfontpath}{fonts/}
\usepackage[book]{islamicbooklet}
% \usepackage[book,text=libron]{islamicbooklet}   % Libron instead of EB Garamond

% Optional page furniture — leave out for none:
% \bookletdomain{www.example.org}    % printed at the foot of the cover
% \coverband{laurel}                 % lighter, more open frame. Default {vine}.

\begin{document}

% ---------------------------------------------------------------------
%  THE COVER. Every key is optional — delete the ones you do not need
%  and the space they occupied closes up. Omit artitle and the whole
%  Arabic block disappears, which is the English-only cover.
%
%  ⚠ A VALUE CONTAINING A COMMA MUST BE BRACED, or keyval reads the
%    comma as the end of the key:  credit = {Surname, Firstname -- ...},
%
%  The frame: put your own full-page artwork in the main folder as
%  cover-frame.pdf (or .png) and it is used; without one, a frame is
%  drawn from ornaments.
% ---------------------------------------------------------------------
\cover{
  entitle     = The Book's Title,
  engloss     = The English Meaning,
  enauthor    = Authored by ...,
  enhonorific = ...,
  endates     = (0000-0000H.),
  credit      = {..., ... -- ...},
  edition     = 1\textsuperscript{st} Edition -- Month Year}
% Arabic cover lines take the keys artitle, arauthor, arhonorific,
% ardates, areditor — type the Arabic yourself.

% The running head, and the page furniture for everything after it.
\bookstyle{The Book's Title}

% Front matter.
\translitchart
\clearpage   % the chart owns its page
\contents

\clearpage

% ---------------------------------------------------------------------
%  THE BODY.
%
%  \heading{...} is numbered automatically and takes English only.
%  Arabic goes in an \ayah, \hadith or \athar block, where it gets the
%  right face and colour.
%
%  For "40th CHAPTER: ..." instead of a bare "40.", put
%  \headinglabel{CHAPTER} in the preamble.
% ---------------------------------------------------------------------
\heading{The First Chapter}

Ordinary prose. Paragraphs are flush left with a small gap between them
and no first-line indent. Just leave a blank line for a new paragraph.

\begin{ayah}
  \framing{Allah said:}
  % \ar{...}  <- UNCOMMENT and type the Arabic here.
  %    Left commented so this starter builds as shipped: the Arabic
  %    faces carry no Latin glyphs, so a placeholder like "..." is a
  %    fatal "there are no glyphs" error, not a missing character.
  \en{...}
  \surah{Al-Fātiḥah}{1:5}
\end{ayah}

\begin{hadith}
  \framing{The Messenger of Allah \saw{} said:}
  % \ar{...}  <- UNCOMMENT and type the Arabic here.
  \en{...}
  \src{Al-Bukhārī, no. 0000}
\end{hadith}

\begin{athar}
  \framing{Ibn ʿAbbās \ra{} said:}
  \en{...}
  \src{...}
\end{athar}

\divider

\heading{The Second Chapter}

Use \tr{transliteration} for Arabic words written in Latin letters,
\g{a grey round gloss} inside a quotation, \ed{a grey square editorial
note}, and \hon{may Allah be pleased with him} for an honorific written
out in words. \link{https://example.com}{A live link} works in a
footnote too.\footnote{\bareurl{https://example.com}}

\end{document}
